\chapter*{Prólogo}
\markboth{Prólogo}{Prólogo}

\section*{¿Por qué este libro?}
Cuando un estudiante de Ciencia de la Computación se enfrenta por primera vez a las matemáticas discretas, suele hacer dos preguntas fundamentales:
\begin{itemize}
  \item ¿para qué sirve esto?
  \item ¿cómo se relaciona con la programación que ya conozco?
\end{itemize}
Este libro es mi intento de responder a esas preguntas de la manera más clara y práctica posible.

\section*{Estructura del libro}
Este libro es la continuación natural de \emph{Estructuras Discretas I} y está organizado en cinco partes principales:
\begin{enumerate}
  \item \textbf{Combinatoria}: El arte del conteo. Desde los principios básicos hasta las aplicaciones en el análisis de algoritmos.
  \item \textbf{Teoría de Grafos}: El estudio de redes y relaciones. Desde las definiciones fundamentales hasta los algoritmos de optimización.
  \item \textbf{Criptografía}: La ciencia de la comunicación segura. Desde los cifrados clásicos hasta el sistema RSA.
  \item \textbf{Programación Funcional con Haskell}: Una introducción práctica a Haskell como herramienta para implementar y verificar conceptos matemáticos.
  \item \textbf{Guías de Práctica}: Ejercicios y proyectos que integran todos los conceptos aprendidos.
\end{enumerate}

\section*{¿Por qué Haskell?}
Haskell no es solo otro lenguaje de programación. Es una herramienta que permite:
\begin{itemize}
  \item Escribir funciones que son matemáticamente puras.
  \item Razonar sobre el código mediante demostraciones formales.
  \item Implementar estructuras de datos con una elegancia que refleja su definición matemática.
  \item Verificar experimentalmente los conceptos teóricos.
\end{itemize}

\section*{Agradecimientos}
A los estudiantes de la Escuela Profesional de Ciencia de la Computación, que con sus preguntas y entusiasmo han enriquecido este material.

\vspace{1cm}
\begin{flushright}
Wilber Roberto Ramos Lovón\\
Arequipa, 2026
\end{flushright}